% =========================================================
% Composition, Inverses, and Image / Preimage Laws
% =========================================================

\subsection{Composition, Inverses, and Set-Image Laws}

% ---------------------------------------------------------
% TOOLKIT
% ---------------------------------------------------------

\begin{definitionbox}{Definition (Composition)}
\begin{definition}[Composition]\label{def:composition}
Let $f : A \to B$ and $g : B \to C$. The \emph{composition} $g \circ f : A \to C$ is:
\[
(g \circ f)(a) := g\bigl(f(a)\bigr)
\quad\text{for all } a \in A.
\]
\end{definition}
\end{definitionbox}

\begin{remark*}[Standard quantified statement]
$\forall a\in A,\; (g\circ f)(a)=g(f(a))$.
\end{remark*}

\begin{remark*}[Interpretation]
Composition applies $f$ first and then applies $g$ to the output. The codomain
of $f$ must match the domain of $g$ so that the second evaluation is defined.
\end{remark*}

\begin{remark*}[Exposition]
Composition is generally not commutative: $g \circ f \neq f \circ g$ even
when both are defined.
\end{remark*}

\begin{dependencies}
\begin{itemize}
  \item \hyperref[def:function]{Function}
  \item \hyperref[def:cartesian-product-rel]{Cartesian Product}
\end{itemize}
\end{dependencies}

\begin{theorembox}{Theorem (Associativity of Composition)}
\begin{theorem}[Associativity of Composition]\label{thm:comp-assoc}
For $f : A \to B$, $g : B \to C$, $h : C \to D$:
\[
h \circ (g \circ f) = (h \circ g) \circ f.
\]
\hyperref[prf:comp-assoc]{Go to proof.}
\end{theorem}
\end{theorembox}

\begin{remark*}[Standard quantified statement]
$\forall f\!:\!A\!\to\!B,\; g\!:\!B\!\to\!C,\; h\!:\!C\!\to\!D,\; \forall a \in A:\; [h \circ (g \circ f)](a) = [(h \circ g) \circ f](a)$.
\end{remark*}

\begin{remark*}[Predicate reading]
$\operatorname{AssociativeComposition}(f,g,h)$ means that the two possible
parenthesizations of $h\circ g\circ f$ agree as functions on $A$.
\end{remark*}

\begin{remark*}[Interpretation]
Associativity says parentheses do not matter for a composable string of three
functions. Only the order of the functions matters.
\end{remark*}

\begin{dependencies}
\begin{itemize}
  \item \hyperref[def:function]{Function}
  \item \hyperref[def:composition]{Composition}
\end{itemize}
\end{dependencies}

\begin{theorembox}{Theorem (Identity and Composition)}
\begin{theorem}[Identity and Composition]\label{thm:comp-id}
For any $f : A \to B$:
\[
f \circ \id_A = f
\quad\text{and}\quad
\id_B \circ f = f.
\]
\hyperref[prf:comp-id]{Go to proof.}
\end{theorem}
\end{theorembox}

\begin{remark*}[Standard quantified statement]
$\forall f : A \to B,\; \forall a \in A:\; (f \circ \id_A)(a) = f(a)$ and $(\id_B \circ f)(a) = f(a)$.
\end{remark*}

\begin{remark*}[Predicate reading]
$\operatorname{IdentityNeutralForComposition}(f,\id_A,\id_B)$ means that the
identity on the domain is a right identity for $f$ and the identity on the
codomain is a left identity for $f$.
\end{remark*}

\begin{remark*}[Interpretation]
Identity functions do nothing under composition. They allow the function $f$ to
be compared with longer composites without changing its action.
\end{remark*}

\begin{dependencies}
\begin{itemize}
  \item \hyperref[def:function]{Function}
  \item \hyperref[def:composition]{Composition}
  \item \hyperref[def:identity]{Identity Function}
\end{itemize}
\end{dependencies}

\begin{theorembox}{Theorem (Injectivity and Surjectivity Under Composition)}
\begin{theorem}[Injectivity and Surjectivity Under Composition]\label{thm:comp-inj-surj}
Let $f : A \to B$ and $g : B \to C$.
\begin{enumerate}[label=(\roman*)]
\item If $f$ and $g$ are injective, then $g \circ f$ is injective.
\item If $f$ and $g$ are surjective, then $g \circ f$ is surjective.
\item If $g \circ f$ is injective, then $f$ is injective.
\item If $g \circ f$ is surjective, then $g$ is surjective.
\end{enumerate}
\hyperref[prf:comp-inj-surj]{Go to proof.}
\end{theorem}
\end{theorembox}

\begin{remark*}[Standard quantified statement]
(i) $(\forall a_1,a_2 \in A,\; f(a_1)=f(a_2)\to a_1=a_2) \wedge (\forall b_1,b_2\in B,\; g(b_1)=g(b_2)\to b_1=b_2) \;\to\; \forall a_1,a_2\in A,\; g(f(a_1))=g(f(a_2))\to a_1=a_2$.\\
(iii) $(\forall a_1,a_2\in A,\; g(f(a_1))=g(f(a_2))\to a_1=a_2) \;\to\; \forall a_1,a_2\in A,\; f(a_1)=f(a_2)\to a_1=a_2$.
\end{remark*}

\begin{remark*}[Predicate reading]
$\operatorname{Injective}(f)$ and $\operatorname{Injective}(g)$ imply
$\operatorname{Injective}(g\circ f)$. The predicate
$\operatorname{Surjective}(f)$ and $\operatorname{Surjective}(g)$ imply
$\operatorname{Surjective}(g\circ f)$. Conversely,
$\operatorname{Injective}(g\circ f)$ implies $\operatorname{Injective}(f)$, and
$\operatorname{Surjective}(g\circ f)$ implies $\operatorname{Surjective}(g)$.
\end{remark*}

\begin{remark*}[Interpretation]
Composing two injective maps preserves the absence of collisions, and composing
two surjective maps preserves the ability to hit every target. If a composite
has one of these properties, only the factor closest to the relevant side is
forced to have it.
\end{remark*}

\begin{remark*}[Exposition]
In (iii), $g$ need not be injective. In (iv), $f$ need not be surjective.
Bijectivity of $g \circ f$ does not imply bijectivity of either $f$ or $g$
individually.
\end{remark*}

\begin{dependencies}
\begin{itemize}
  \item \hyperref[def:function]{Function}
  \item \hyperref[def:composition]{Composition}
  \item \hyperref[def:injective]{Injective Function}
  \item \hyperref[def:surjective]{Surjective Function}
\end{itemize}
\end{dependencies}

% ---------------------------------------------------------
% Inverse functions
% ---------------------------------------------------------
\begin{definitionbox}{Definition (Invertible Function)}
\begin{definition}[Invertible Function]\label{def:invertible-function}
Let $f : A \to B$ be a function. We say that $f$ is \emph{invertible} if
there exists a function $g : B \to A$ such that
\[
g \circ f = \id_A
\quad\text{and}\quad
f \circ g = \id_B.
\]
Such a function $g$ is called an \emph{inverse} of $f$.
\end{definition}
\end{definitionbox}

\begin{remark*}[Standard quantified statement]
$\operatorname{Invertible}(f)\;\leftrightarrow\;\exists g:B\to A,\;
g\circ f=\id_A \wedge f\circ g=\id_B$.
\end{remark*}

\begin{remark*}[Interpretation]
An invertible function can be undone from both sides. The same inverse returns
inputs after applying $f$ and also reaches each codomain element from its chosen
input.
\end{remark*}

\begin{dependencies}
\begin{itemize}
  \item \hyperref[def:function]{Function}
  \item \hyperref[def:composition]{Composition}
  \item \hyperref[def:identity]{Identity Function}
\end{itemize}
\end{dependencies}

\begin{definitionbox}{Definition (Inverse Function)}
\begin{definition}[Inverse Function]\label{def:inverse}
Let $f : A \to B$ be bijective. The \emph{inverse function} is
$f^{-1} : B \to A$ defined by: $f^{-1}(b) = a$ iff $f(a) = b$.
\end{definition}
\end{definitionbox}

\begin{remark*}[Standard quantified statement]
$\forall b\in B,\; f^{-1}(b)=a \leftrightarrow f(a)=b$.
\end{remark*}

\begin{remark*}[Interpretation]
The inverse function reverses a bijection by sending each output back to its
unique input.
\end{remark*}

\begin{dependencies}
\begin{itemize}
  \item \hyperref[def:function]{Function}
  \item \hyperref[def:invertible-function]{Invertible Function}
  \item \hyperref[def:bijective]{Bijective Function}
  \item \hyperref[def:identity]{Identity Function}
\end{itemize}
\end{dependencies}

\begin{theorembox}{Theorem (Characterization of Inverse Functions)}
\begin{theorem}[Characterization of Inverse Functions]\label{thm:inverse-char}
Let $f : A \to B$ be bijective. Then
\[
f^{-1} \circ f = \id_A
\quad\text{and}\quad
f \circ f^{-1} = \id_B.
\]
Conversely, a function admits an inverse iff it is bijective.
\hyperref[prf:inverse-char]{Go to proof.}
\end{theorem}
\end{theorembox}

\begin{remark*}[Standard quantified statement]
$\forall a \in A:\; f^{-1}(f(a)) = a$ and $\forall b \in B:\; f(f^{-1}(b)) = b$.\\
Converse: $(\exists g: B\to A,\; g\circ f = \id_A \wedge f\circ g = \id_B) \;\leftrightarrow\; f$ is bijective.
\end{remark*}

\begin{remark*}[Predicate reading]
$\operatorname{Bijective}(f)$ is equivalent to the existence of a function
$g:B\to A$ satisfying $\operatorname{LeftInverse}(g,f)$ and
$\operatorname{RightInverse}(g,f)$.
\end{remark*}

\begin{remark*}[Interpretation]
Being bijective is exactly the condition that makes the inverse function a
two-sided inverse. Injectivity supplies uniqueness of the input, and
surjectivity supplies existence of an input.
\end{remark*}

\begin{dependencies}
\begin{itemize}
  \item \hyperref[def:function]{Function}
  \item \hyperref[def:composition]{Composition}
  \item \hyperref[def:invertible-function]{Invertible Function}
  \item \hyperref[def:inverse]{Inverse Function}
  \item \hyperref[def:bijective]{Bijective Function}
  \item \hyperref[def:identity]{Identity Function}
\end{itemize}
\end{dependencies}

\begin{theorembox}{Theorem (Inverse of a Composition)}
\begin{theorem}[Inverse of a Composition]\label{thm:inverse-comp}
Let $f : A \to B$ and $g : B \to C$ be bijective. Then $g \circ f$ is
bijective and
\[
(g \circ f)^{-1} = f^{-1} \circ g^{-1}.
\]
\hyperref[prf:inverse-comp]{Go to proof.}
\end{theorem}
\end{theorembox}

\begin{remark*}[Standard quantified statement]
$\forall c \in C:\; (g \circ f)^{-1}(c) = f^{-1}(g^{-1}(c))$.
\end{remark*}

\begin{remark*}[Interpretation]
To undo a composite, undo the last function first. The inverse of a composition
therefore reverses the order of the original functions.
\end{remark*}

\begin{dependencies}
\begin{itemize}
  \item \hyperref[def:function]{Function}
  \item \hyperref[def:composition]{Composition}
  \item \hyperref[def:inverse]{Inverse Function}
  \item \hyperref[def:bijective]{Bijective Function}
\end{itemize}
\end{dependencies}

% ---------------------------------------------------------
% One-sided inverses
% ---------------------------------------------------------
\begin{definitionbox}{Definition (Left Inverse)}
\begin{definition}[Left Inverse]\label{def:left-inverse-function}
Let $f : A \to B$.
A \emph{left inverse} of $f$ is a function $g : B \to A$ with $g \circ f = \id_A$.
\end{definition}
\end{definitionbox}

\begin{remark*}[Standard quantified statement]
$\operatorname{LeftInverse}(g,f)\;\leftrightarrow\; g:B\to A \wedge
g\circ f=\id_A$.
\end{remark*}

\begin{remark*}[Interpretation]
A left inverse recovers every input after applying $f$. It rules out collisions
among inputs of $f$.
\end{remark*}

\begin{dependencies}
\begin{itemize}
  \item \hyperref[def:function]{Function}
  \item \hyperref[def:composition]{Composition}
  \item \hyperref[def:identity]{Identity Function}
\end{itemize}
\end{dependencies}

\begin{definitionbox}{Definition (Right Inverse)}
\begin{definition}[Right Inverse]\label{def:right-inverse-function}
Let $f : A \to B$.
A \emph{right inverse} (or \emph{section}) of $f$ is a function $h : B \to A$
with $f \circ h = \id_B$.
\end{definition}
\end{definitionbox}

\begin{remark*}[Standard quantified statement]
$\operatorname{RightInverse}(h,f)\;\leftrightarrow\; h:B\to A \wedge
f\circ h=\id_B$.
\end{remark*}

\begin{remark*}[Interpretation]
A right inverse chooses an input above each codomain element. It can exist only
when every codomain element is hit by $f$.
\end{remark*}

\begin{remark*}[Exposition]
Left inverses and right inverses are the two one-sided inverse notions.
A left inverse reverses $f$ on the domain side, so it can exist only when
different inputs do not collapse to the same output. A right inverse chooses
one input over each codomain element, so it can exist only when every codomain
element is actually hit.
\end{remark*}

\begin{dependencies}
\begin{itemize}
  \item \hyperref[def:function]{Function}
  \item \hyperref[def:composition]{Composition}
  \item \hyperref[def:identity]{Identity Function}
\end{itemize}
\end{dependencies}

\begin{figure}[htbp]
\centering
\begin{tikzpicture}[atlas, x=1cm, y=1cm, >=Stealth, every node/.style={font=\small}]
  \node[atlas label] at (-2.0,2.65) {Left inverse};
  \draw[atlasgray, line width=0.7pt] (-3.75,-0.15) ellipse (0.65 and 1.65);
  \draw[atlasgray, line width=0.7pt] (-1.55,-0.15) ellipse (0.75 and 1.65);
  \node[atlas label] at (-3.75,1.75) {$A$};
  \node[atlas label] at (-1.55,1.75) {$B$};

  \node[atlas dot, label={[atlas label,left]$a_1$}] (la1) at (-3.75,0.65) {};
  \node[atlas dot, label={[atlas label,left]$a_2$}] (la2) at (-3.75,-0.95) {};
  \node[atlas dot, label={[atlas label,right]$f(a_1)$}] (lb1) at (-1.55,0.65) {};
  \node[atlas dot, label={[atlas label,right]$f(a_2)$}] (lb2) at (-1.55,-0.95) {};
  \node[atlas dot, label={[atlas label,right]$b$}] (lb3) at (-1.55,-0.15) {};

  \draw[atlasblue, atlas probe, ->] (la1) -- (lb1);
  \draw[atlasblue, atlas probe, ->] (la2) -- (lb2);
  \draw[atlasgreen, atlas probe, ->, bend left=18] (lb1) to (la1);
  \draw[atlasgreen, atlas probe, ->, bend right=18] (lb2) to (la2);
  \draw[atlasgray, line width=0.8pt, dashed, ->] (lb3) to[bend left=8] (la1);
  \node[atlas label] at (-2.65,-2.0) {$g\circ f=\id_A$};

  \node[atlas label] at (2.0,2.65) {Right inverse};
  \draw[atlasgray, line width=0.7pt] (0.25,-0.15) ellipse (0.75 and 1.65);
  \draw[atlasgray, line width=0.7pt] (2.45,-0.15) ellipse (0.65 and 1.65);
  \node[atlas label] at (0.25,1.75) {$A$};
  \node[atlas label] at (2.45,1.75) {$B$};

  \node[atlas dot, label={[atlas label,left]$a_1$}] (ra1) at (0.25,0.85) {};
  \node[atlas dot, label={[atlas label,left]$a_2$}] (ra2) at (0.25,0.05) {};
  \node[atlas dot, label={[atlas label,left]$a_3$}] (ra3) at (0.25,-1.05) {};
  \node[atlas dot, label={[atlas label,right]$b_1$}] (rb1) at (2.45,0.45) {};
  \node[atlas dot, label={[atlas label,right]$b_2$}] (rb2) at (2.45,-1.05) {};

  \draw[atlasblue, atlas probe, ->] (ra1) -- (rb1);
  \draw[atlasblue, atlas probe, ->] (ra2) -- (rb1);
  \draw[atlasblue, atlas probe, ->] (ra3) -- (rb2);
  \draw[atlasgreen, atlas probe, ->, bend left=18] (rb1) to (ra1);
  \draw[atlasgreen, atlas probe, ->, bend right=18] (rb2) to (ra3);
  \node[atlas label] at (1.35,-2.0) {$f\circ h=\id_B$};
\end{tikzpicture}

\caption{Left and right inverses impose different one-sided requirements.}
\label{fig:one-sided-inverses}
\end{figure}

\begin{theorembox}{Theorem (One-Sided Inverses and Function Properties)}
\begin{theorem}[One-Sided Inverses and Function Properties]\label{thm:one-sided}
Let $f : A \to B$.
\begin{enumerate}[label=(\roman*)]
\item $f$ has a left inverse $\iff$ $f$ is injective (and $A \neq \varnothing$ or $A=B=\varnothing$).
\item $f$ has a right inverse $\iff$ $f$ is surjective.
\item If $f$ has both a left inverse $g$ and a right inverse $h$, then $g = h$
and $f$ is bijective.
\end{enumerate}
\hyperref[prf:one-sided]{Go to proof.}
\end{theorem}
\end{theorembox}

\begin{remark*}[Standard quantified statement]
(i) $(\exists g: B\to A,\; \forall a\in A,\; g(f(a))=a) \;\leftrightarrow\; \forall a_1,a_2\in A,\; f(a_1)=f(a_2)\to a_1=a_2$.\\
(ii) $(\exists h: B\to A,\; \forall b\in B,\; f(h(b))=b) \;\leftrightarrow\; \forall b\in B,\; \exists a\in A,\; f(a)=b$.
\end{remark*}

\begin{remark*}[Predicate reading]
$\exists g\,\operatorname{LeftInverse}(g,f)$ is equivalent to
$\operatorname{Injective}(f)$, with the stated empty-domain convention.
$\exists h\,\operatorname{RightInverse}(h,f)$ is equivalent to
$\operatorname{Surjective}(f)$. Having both one-sided inverses forces
$\operatorname{Bijective}(f)$.
\end{remark*}

\begin{remark*}[Interpretation]
Left inverses detect whether information was lost, while right inverses detect
whether every target can be chosen from. When both exist, the two choices agree
and give the ordinary inverse.
\end{remark*}

\begin{remark*}[Exposition]
The existence of a right inverse for every surjective function is equivalent
to the Axiom of Choice.
\end{remark*}

\begin{dependencies}
\begin{itemize}
  \item \hyperref[def:function]{Function}
  \item \hyperref[def:composition]{Composition}
  \item \hyperref[def:left-inverse-function]{Left Inverse}
  \item \hyperref[def:right-inverse-function]{Right Inverse}
  \item \hyperref[def:injective]{Injective Function}
  \item \hyperref[def:surjective]{Surjective Function}
  \item \hyperref[def:bijective]{Bijective Function}
\end{itemize}
\end{dependencies}

% ---------------------------------------------------------
% Restriction and Extension
% ---------------------------------------------------------
\begin{definitionbox}{Definition (Restriction)}
\begin{definition}[Restriction]\label{def:restriction}
Let $f : A \to B$ and $C \subseteq A$. The \emph{restriction} of $f$ to $C$
is $f|_C : C \to B$, $f|_C(a) = f(a)$ for all $a \in C$.
\end{definition}
\end{definitionbox}

\begin{remark*}[Standard quantified statement]
$\forall a\in C,\; f|_C(a)=f(a)$.
\end{remark*}

\begin{remark*}[Interpretation]
Restriction keeps the same rule but narrows the domain. This can remove inputs
that caused unwanted behavior, such as collisions.
\end{remark*}

\begin{dependencies}
\begin{itemize}
  \item \hyperref[def:function]{Function}
  \item \hyperref[def:subset]{Subset}
  \item \hyperref[def:domain]{Domain}
\end{itemize}
\end{dependencies}

\begin{definitionbox}{Definition (Extension)}
\begin{definition}[Extension]\label{def:fn-extension}
Let $A \subseteq A'$, and let $f : A \to B$. A function $g : A' \to B$ is an
\emph{extension} of $f$ if $g|_A = f$.
\end{definition}
\end{definitionbox}

\begin{remark*}[Standard quantified statement]
$\operatorname{Extension}(g,f)\;\leftrightarrow\; A\subseteq A' \wedge
g:A'\to B \wedge \forall a\in A,\; g(a)=f(a)$.
\end{remark*}

\begin{remark*}[Interpretation]
An extension preserves the old function on its original domain while assigning
values on a larger domain.
\end{remark*}

\begin{remark*}[Exposition]
An extension agrees with $f$ on all of $A$ but may be defined on a larger
domain $A'$. Extensions are generally not unique; uniqueness requires additional
constraints (such as continuity in analysis or linearity in algebra), leading to
important results like the Tietze Extension Theorem and the Hahn--Banach
Theorem.
\end{remark*}

\begin{remark*}[Exposition]
The chapter \hyperref[ch:functions-and-order]{Functions and Order} studies
how restricting the domain, corestricting to the image, or enlarging the
codomain changes injectivity, surjectivity, inverse-on-image behavior, and
monotonicity.
\end{remark*}

\begin{dependencies}
\begin{itemize}
  \item \hyperref[def:function]{Function}
  \item \hyperref[def:subset]{Subset}
  \item \hyperref[def:restriction]{Restriction}
  \item \hyperref[def:domain]{Domain}
  \item \hyperref[def:codomain]{Codomain}
\end{itemize}
\end{dependencies}

% ---------------------------------------------------------
% Image and Preimage Laws
% ---------------------------------------------------------
\begin{theorembox}{Theorem (Preimages Preserve Set Operations)}
\begin{theorem}[Preimages Preserve Set Operations]\label{thm:preimage-ops}
Let $f : A \to B$ and $S, T \subseteq B$. Then
\begin{enumerate}[label=(\roman*)]
\item $f^{-1}(S \cup T) = f^{-1}(S) \cup f^{-1}(T)$,
\item $f^{-1}(S \cap T) = f^{-1}(S) \cap f^{-1}(T)$,
\item $f^{-1}(S \setminus T) = f^{-1}(S) \setminus f^{-1}(T)$,
\item $f^{-1}(S^c) = (f^{-1}(S))^c$.
\end{enumerate}
\hyperref[prf:preimage-ops]{Go to proof.}
\end{theorem}
\end{theorembox}

\begin{remark*}[Standard quantified statement]
(i) $\forall a \in A:\; a \in f^{-1}(S\cup T) \leftrightarrow (f(a)\in S \vee f(a)\in T)$.\\
(ii) $\forall a \in A:\; a \in f^{-1}(S\cap T) \leftrightarrow (f(a)\in S \wedge f(a)\in T)$.\\
(iv) $\forall a \in A:\; a \in f^{-1}(S^c) \leftrightarrow f(a)\notin S$.
\end{remark*}

\begin{remark*}[Predicate reading]
$\operatorname{PreservesUnion}(f,S,T)$,
$\operatorname{PreservesIntersection}(f,S,T)$,
$\operatorname{PreservesDifference}(f,S,T)$, and
$\operatorname{PreservesComplement}(f,S)$ mean that preimage converts
the corresponding set operation in $B$ into the same set operation in $A$.
\end{remark*}

\begin{remark*}[Interpretation]
Preimage laws are exact because membership in a preimage is just a statement
about where $f(a)$ belongs. Logical connectives are preserved when pulled back
to the domain.
\end{remark*}

\begin{dependencies}
\begin{itemize}
  \item \hyperref[def:function]{Function}
  \item \hyperref[def:preimage]{Preimage}
  \item \hyperref[def:union]{Union}
  \item \hyperref[def:intersection]{Intersection}
  \item \hyperref[def:set-difference]{Set Difference}
  \item \hyperref[def:complement]{Complement}
\end{itemize}
\end{dependencies}

\subsubsection*{Indexed Preimage Operations}

\begin{theorembox}{Theorem (Preimages Preserve Indexed Set Operations)}
\begin{theorem}[Preimages Preserve Indexed Set Operations]\label{thm:indexed-preimage-ops}
Let $f:A\to B$ and let $\{T_i\}_{i\in I}$ be an indexed family of subsets of
$B$. Then
\[
f^{-1}\!\left(\bigcup_{i\in I}T_i\right)
=
\bigcup_{i\in I}f^{-1}(T_i).
\]
If $I\neq\varnothing$, then
\[
f^{-1}\!\left(\bigcap_{i\in I}T_i\right)
=
\bigcap_{i\in I}f^{-1}(T_i).
\]
\hyperref[prf:indexed-preimage-ops]{Go to proof.}
\end{theorem}
\end{theorembox}

\begin{remark*}[Standard quantified statement]
$\forall a\in A:\; a\in f^{-1}(\bigcup_{i\in I}T_i)
\leftrightarrow \exists i\in I,\; f(a)\in T_i$.\\
If $I\neq\varnothing$, then
$\forall a\in A:\; a\in f^{-1}(\bigcap_{i\in I}T_i)
\leftrightarrow \forall i\in I,\; f(a)\in T_i$.
\end{remark*}

\begin{remark*}[Predicate reading]
$\operatorname{PreservesIndexedUnion}(f,\{T_i\}_{i\in I})$ and
$\operatorname{PreservesIndexedIntersection}(f,\{T_i\}_{i\in I})$ mean
that preimage commutes with the indexed union, and with the indexed intersection
when the index set is nonempty.
\end{remark*}

\begin{remark*}[Interpretation]
Preimage turns membership in a union or intersection back into the same
logical condition on the input side. When $I=\mathbb{N}$, these identities say
that preimages commute with countable unions and countable intersections.
Topology uses the arbitrary union case for open sets and the finite, countable,
or arbitrary intersection cases according to the structure under discussion.
\end{remark*}

\begin{dependencies}
\begin{itemize}
  \item \hyperref[def:function]{Function}
  \item \hyperref[def:preimage]{Preimage}
  \item \hyperref[def:indexed-family]{Indexed Family of Sets}
  \item \hyperref[def:indexed-union]{Indexed Union}
  \item \hyperref[def:indexed-intersection]{Indexed Intersection}
\end{itemize}
\end{dependencies}

\begin{theorembox}{Theorem (Images and Set Operations)}
\begin{theorem}[Images and Set Operations]\label{thm:image-ops}
Let $f : A \to B$ and $S, T \subseteq A$. Then
\begin{enumerate}[label=(\roman*)]
\item $f(S \cup T) = f(S) \cup f(T)$,
\item $f(S \cap T) \subseteq f(S) \cap f(T)$,
\item $f(S \setminus T) \supseteq f(S) \setminus f(T)$.
\end{enumerate}
Equality holds in (ii) and (iii) for all $S,T$ iff $f$ is injective.
\hyperref[prf:image-ops]{Go to proof.}
\end{theorem}
\end{theorembox}

\begin{remark*}[Standard quantified statement]
(i) $\forall b\in B:\; b \in f(S\cup T) \leftrightarrow (\exists a\in S,\; f(a)=b) \vee (\exists a\in T,\; f(a)=b)$.\\
(ii) $\forall b\in B:\; (\exists a\in S\cap T,\; f(a)=b) \to (\exists a\in S,\; f(a)=b) \wedge (\exists a\in T,\; f(a)=b)$; equality iff $f$ injective.
\end{remark*}

\begin{remark*}[Predicate reading]
$\operatorname{PreservesUnion}(f,S,T)$ means that $f(S\cup T)=f(S)\cup
f(T)$. The predicates $\operatorname{MapsInto}(f,S,T)$ and
$\operatorname{MapsInto}(f,S,T)$ record the one-sided image
laws, with equality uniformly when $\operatorname{Injective}(f)$ holds.
\end{remark*}

\begin{remark*}[Interpretation]
Forward images preserve unions because an input can come from either side.
Intersections and differences lose information unless the function is
injective.
\end{remark*}

\begin{dependencies}
\begin{itemize}
  \item \hyperref[def:function]{Function}
  \item \hyperref[def:image-set]{Image of a Set}
  \item \hyperref[def:union]{Union}
  \item \hyperref[def:intersection]{Intersection}
  \item \hyperref[def:set-difference]{Set Difference}
  \item \hyperref[def:injective]{Injective Function}
\end{itemize}
\end{dependencies}

\begin{theorembox}{Theorem (Image--Preimage Adjunction)}
\begin{theorem}[Image--Preimage Adjunction]\label{thm:image-preimage-adjunction}
Let $f:A\to B$, let $S\subseteq A$, and let $T\subseteq B$. Then
\[
S\subseteq f^{-1}(f(S)),
\qquad
f(f^{-1}(T))\subseteq T.
\]
Moreover, $S=f^{-1}(f(S))$ for every $S\subseteq A$ iff $f$ is injective, and
$f(f^{-1}(T))=T$ for every $T\subseteq B$ iff $f$ is surjective.
\hyperref[prf:image-preimage-adjunction]{Go to proof.}
\end{theorem}
\end{theorembox}

\begin{remark*}[Standard quantified statement]
$\forall S\subseteq A,\; S\subseteq f^{-1}(f(S))$ and
$\forall T\subseteq B,\; f(f^{-1}(T))\subseteq T$. Moreover,
$(\forall S\subseteq A,\; S=f^{-1}(f(S)))\leftrightarrow
\operatorname{Injective}(f)$, and $(\forall T\subseteq B,\;
f(f^{-1}(T))=T)\leftrightarrow \operatorname{Surjective}(f)$.
\end{remark*}

\begin{remark*}[Predicate reading]
$\operatorname{ImagePreimageAdjunction}(f,S,T)$ records the inclusions
$S\subseteq f^{-1}(f(S))$ and $f(f^{-1}(T))\subseteq T$, together with the
injective and surjective equality criteria.
\end{remark*}

\begin{remark*}[Interpretation]
Sending a set forward and then pulling it back can add points that collide with
the original set. Pulling back and then sending forward cannot leave the target
set. Injectivity and surjectivity are exactly the conditions that turn these
inclusions into equalities.
\end{remark*}

\begin{dependencies}
\begin{itemize}
  \item \hyperref[def:function]{Function}
  \item \hyperref[def:image-set]{Image of a Set}
  \item \hyperref[def:preimage]{Preimage}
  \item \hyperref[def:injective]{Injective Function}
  \item \hyperref[def:surjective]{Surjective Function}
\end{itemize}
\end{dependencies}

\begin{propositionbox}{Proposition (Preimage Under an Inverse Function)}
\begin{proposition}[Preimage Under an Inverse Function]\label{prop:preimage-under-inverse-function}
Let $f:A\to B$ be bijective, and let $S\subseteq A$. Then
\[
(f^{-1})^{-1}(S)=f(S),
\]
where $(f^{-1})^{-1}(S)$ denotes the preimage of $S$ under the inverse
function $f^{-1}:B\to A$.
\hyperref[prf:preimage-under-inverse-function]{Go to proof.}
\end{proposition}
\end{propositionbox}

\begin{remark*}[Standard quantified statement]
$\forall b\in B,\; b\in (f^{-1})^{-1}(S)\leftrightarrow b\in f(S)$.
\end{remark*}

\begin{remark*}[Interpretation]
Preimage under the inverse function is the same operation as forward image
under the original bijection.
\end{remark*}

\begin{remark*}[Exposition]
The inner $f^{-1}$ is the inverse function. The outer ${}^{-1}(S)$ is
preimage notation applied to that inverse function. Thus the proposition says
that pulling $S\subseteq A$ back along $f^{-1}:B\to A$ gives the same subset of
$B$ as pushing $S$ forward along $f:A\to B$.
\end{remark*}

\begin{remark*}[Exposition]
\mbox{}\par\smallskip
\small
\begin{tabularx}{\textwidth}{l X X}
\toprule
\textbf{Operation} & \textbf{Preimage law, for $S,T\subseteq B$} & \textbf{Image law, for $C,D\subseteq A$} \\
\midrule
Union
& $f^{-1}(S\cup T)=f^{-1}(S)\cup f^{-1}(T)$
& $f(C\cup D)=f(C)\cup f(D)$ \\
Intersection
& $f^{-1}(S\cap T)=f^{-1}(S)\cap f^{-1}(T)$
& $f(C\cap D)\subseteq f(C)\cap f(D)$; equality for all $C,D$ iff $f$ is injective \\
Difference
& $f^{-1}(S\setminus T)=f^{-1}(S)\setminus f^{-1}(T)$
& $f(C\setminus D)\supseteq f(C)\setminus f(D)$; equality for all $C,D$ iff $f$ is injective \\
Complement
& $f^{-1}(B\setminus S)=A\setminus f^{-1}(S)$
& In general, $f(A\setminus C)$ need not equal $B\setminus f(C)$ \\
\bottomrule
\end{tabularx}
\end{remark*}

\begin{remark*}[Exposition]
\mbox{}\par\smallskip
\small
\begin{tabularx}{\textwidth}{l X X}
\toprule
\textbf{Construction} & \textbf{Set-operation interaction} & \textbf{Reference} \\
\midrule
Composition
& If $f:A\to B$, $g:B\to C$, and $U\subseteq C$, then
$(g\circ f)^{-1}(U)=f^{-1}(g^{-1}(U))$.
& \hyperref[def:composition]{Composition}, \hyperref[def:preimage]{Preimage} \\
Inverse function
& If $f:A\to B$ is bijective, inverse images along $f^{-1}$ are ordinary
preimages for the function $f^{-1}:B\to A$.
& \hyperref[def:inverse]{Inverse Function}, \hyperref[thm:inverse-char]{Inverse Characterization} \\
Inverse of composition
& If $f$ and $g$ are bijective, then
$(g\circ f)^{-1}=f^{-1}\circ g^{-1}$, so inverse order reverses composition.
& \hyperref[thm:inverse-comp]{Inverse of a Composition} \\
One-sided inverse
& A left inverse detects injectivity; a right inverse detects surjectivity.
These properties control when image inclusions become equalities.
& \hyperref[thm:one-sided]{One-Sided Inverses} \\
\bottomrule
\end{tabularx}
\end{remark*}

\begin{remark*}[Exposition]
Preimages preserve union, intersection, difference, and complement exactly.
Forward images preserve unions exactly, but generally do not preserve
intersections, differences, or complements. The inclusions for intersections
and differences become equalities uniformly when the function is injective.
This asymmetry is fundamental and explains why preimages appear more naturally
in topology and measure theory.
\end{remark*}

\begin{dependencies}
\begin{itemize}
  \item \hyperref[def:image-set]{Image of a Set}
  \item \hyperref[def:preimage]{Preimage}
  \item \hyperref[def:inverse]{Inverse Function}
  \item \hyperref[def:bijective]{Bijective Function}
\end{itemize}
\end{dependencies}
